\documentclass[12pt]{article}
\usepackage{xcolor}
\parindent 0px
\pagestyle{empty}
\begin{document}

\textbf{Soal Latihan JavaScript 1} \\

\textbf{Soal 1 - Comment} \\

Tuliskan kode JavaScript yang berisi: 
\begin{quote}
\begin{itemize}
\item Sebuah single-line comment yang menjelaskan bahwa variabel nama 
menyimpan nama pengguna
\item Sebuah multi-line comment yang menjelaskan program 
akan menampilkan sapaan.
\end{itemize}
\end{quote}

\textbf{Soal 2 - Variabel dan Tipe data} \\

Deklarasikan 4 variabel dengan ketentuan: 
\begin{quote}
\begin{enumerate}
        \item \verb+negara+ bertipe string, diisi nama negara favorit.
        \item \verb+penduduk+ bertipe number, diisi perkiraan jumlah penduduk (dalam juta).
        \item \verb+isAsia+ bertipe boolean, diisi nilai sesuai apakah negara tersebut di Asia.
        \item \verb+dataLain+ bertipe null.      
\end{enumerate}
\end{quote}

\textbf{Soal 4 - Operator (Aritmatika + Perbandingan)} \\

Deklarasikan:

\begin{quote}
\textcolor{orange}{let ${a = 20;}$} \\


\textcolor{blue}{let ${b = 7;}$} \\ 


\textcolor{teal}{let ${c = "20";}$}

\end{quote} 

Lakukan operasi berikut dan tulis hasilnya dalam bentuk komentar \textcolor{red}{tanpa menjalankan kode terlebih dahulu}: 
\underline{(tebak dulu, lalu anda bisa cek dengan menjalankan)}:
\begin{quote}
\begin{enumerate}
        \item $a + b$
        \item $a - b$
        \item $a \ast b$
        \item $a \div b$
        \item $a == c$
        \item $a === c$
        \item ${a > b} \;\&\&\; {b > 0}$
        \item $a > b \Vert  b > 10$
\end{enumerate}
\end{quote} 

\textbf{Soal 5 - Campuran (Konversi + Operator)} \\

Buatlah program sederhana dengan langkah berikut:
\begin{quote}
\begin{enumerate}
        \item Buat variabel \verb+harga+ $=$ "2500" (string).
        \item Buat variabel \verb+jumlah+ $=$ "3" (string).
        \item Konversi kedua variabel tersebut menjadi number.
        \item Hitung total harga dengan Operator $\ast$.
        \item Tampilkan hasil totaldan hasil Perbandingan dengan \verb+console.log+  
\end{enumerate}
\end{quote}

\textbf{Soal 6 - (Bonus) Logic sederhana} \\

Buat variabel \verb+nilaiUjian+ $=$ 85 dan \verb+absensi+ $=$ "75" (string).

\begin{quote}
\begin{enumerate}
        \item Konversi \verb+absensi+ ke number.
        \item Gunakan Operator $+$ untuk menjumlah \verb+nilaiUjian+ dengan \verb+absensi+ (setelah di konversi).
        \item Cek apakah hasil penjumlahan lebih besar dari 150 dan \verb+nilaiUjian+ lebih dari 80.
        \item Simpan hasil pengecekan tersebut ke variabel \verb+lulus+
        \item Tampilkan semua nilai(\verb+nilaiUjian+, \verb+absensi+, \verb+total+, \verb++lulus) dengan rapi.
\end{enumerate}
\end{quote}
\end{document} 